\section{How findings are weighed}
\label{sec:F.1}
\textbf{The corroborated licence, which confirms or lifts the emergency licence later and is the only licence the genocide and occupation exceptions have, rests on} a disinterested finding. The kinds of finding are four: a judgment or advisory opinion of the International Court of Justice, or a judgment of a regional court with adversarial process; the finding of a fact-finding mission or commission of inquiry mandated by a body no party to the conflict controls, and meeting the four conditions below; a determination of the Security Council; and a resolution adopted by two-thirds of the General Assembly members present and voting, where the facts it rests on aren't in dispute. Provisional-measures orders don't count, except as one leg of the genocide exception's two-body rule, since the Court issues them on the plausibility of the rights claimed.

A Council authorization counts as corroboration because it needs the permanent members not to block it, and agreement among rivals whose stakes point different ways is about as disinterested as international judgment gets. A Council determination counts the same way, as a vote that corroborates. The Council's silence on a defense claim counts for nothing, because a veto is a stake.

The four kinds of finding are not equal, and the term weighs them by the process behind them. A merits judgment of the International Court of Justice comes out of adversarial argument, a record and published reasons, and is the strongest; a regional court with adversarial process, such as the European Court of Human Rights, counts as a court for the same reason. A commission of inquiry or fact-finding mission publishes reasons and evidence but hears no adversary; its finding counts where no party to the conflict controls the mandating body and the commission meets the four conditions below, and a later merits judgment displaces it. A Council determination and a General Assembly resolution are votes, not findings of fact. Where the facts aren't in dispute and the question is what they amount to, a vote is a finding; where the facts are disputed, it counts only as corroboration of a record made elsewhere. For the genocide exception the ranking becomes a requirement of two kinds, or the judgment alone, and the occupation exception borrows it. Where a party disputes that a body meets the test above, the dispute and the answer to it go on the record beside the finding. A commission mandated by a belligerent fails whatever its procedure. Past that threshold, whether a commission counts turns on four things that don't depend on who mandated it: its evidence and methodology are published; it offered the parties a reply and recorded the answer or the refusal; a later merits judgment displaces it; and its mandate was fixed before the finding it reaches. A body mandated by a political majority can meet all four and one mandated by a court can fail them. The test is run in \ref{sec:F.2} on a commission I would rather not count, the Organisation for the Prohibition of Chemical Weapons' Investigation and Identification Team, which attributes chemical attacks in Syria.

Whether the facts behind a vote are in dispute is itself a finding, and the belligerent will always dispute them. The rule is that of the discount and the burden (Chapter~\ref{sec:22}): a dispute counts where it is substantive on the documented record, not where it is asserted, and the party that controls the evidence bears the burden of showing it is real. Whoever writes the deployment's floor applies the rule in its published translation, and the floor court hears a contest.\footnote{General Assembly majorities are often bloc votes, and the weighting allows for it: a vote is a finding only where the facts aren't in dispute, and corroborates a record made elsewhere otherwise. The two-thirds threshold is not offered as proof of disinterest. It is offered because a two-thirds majority of a body whose blocs point different ways is harder to assemble on a false record than on a true one, and because the alternative, no vote counting at all, leaves Cyprus and Karabakh with nothing but courts that take decades.}
